\documentclass[11pt,a4paper]{article}
\usepackage[utf8]{inputenc}
\usepackage[spanish]{babel}
\usepackage[T1]{fontenc}
\usepackage{hyperref}
\usepackage{geometry}
\geometry{a4paper, margin=1in}
\usepackage{listings}
\usepackage{xcolor}

\definecolor{codegray}{rgb}{0.5,0.5,0.5}
\definecolor{backcolour}{rgb}{0.95,0.95,0.92}

\lstdefinestyle{mystyle}{
    backgroundcolor=\color{backcolour},   
    commentstyle=\color{codegray},
    keywordstyle=\color{blue},
    numberstyle=\tiny\color{codegray},
    stringstyle=\color{orange},
    basicstyle=\ttfamily\footnotesize,
    breakatwhitespace=false,         
    breaklines=true,                 
    captionpos=b,                    
    keepspaces=true,                 
    numbers=left,                    
    numbersep=5pt,                  
    showspaces=false,                
    showstringspaces=false,
    showtabs=false,                  
    tabsize=2
}

\lstset{style=mystyle}

\title{Guía de Uso Rápido de Docker (Equipo de Desarrollo)}
\author{}
\date{}

\begin{document}

\maketitle

Esta guía documenta los pasos necesarios para que cualquier miembro del equipo pueda levantar el entorno completo de la plataforma de comparación de supermercados (Frontend, Backend, Scraper, Base de Datos) de manera estandarizada y sin tener que instalar lenguajes de programación locales.

\section{Dependencias del Entorno e Instalación}

La arquitectura del proyecto requiere la instalación del motor de \textbf{Docker} (para la construcción de contenedores vía \texttt{Dockerfile}) y \textbf{Docker Compose} (para la orquestación y despliegue del stack de microservicios). A continuación, se detallan los procedimientos de instalación y configuración según el sistema operativo anfitrión.

\subsection{Entornos Linux}

\subsubsection{Arch Linux}
En distribuciones basadas en Arch, los paquetes \texttt{docker} y \texttt{docker-compose} se encuentran disponibles en los repositorios oficiales de la distribución. Ejecute los siguientes comandos:
\begin{lstlisting}[language=bash]
sudo pacman -Syu docker docker-compose
# Habilitar e iniciar el demonio systemd de Docker
sudo systemctl enable --now docker.service
# Adicionar el usuario actual al grupo docker (requiere reinicio de sesion o recarga de grupos)
sudo usermod -aG docker $USER
\end{lstlisting}

\subsubsection{Fedora}
Para entornos basados en Fedora, se recomienda utilizar el repositorio oficial proporcionado por Docker Inc.:
\begin{lstlisting}[language=bash]
sudo dnf install dnf-plugins-core
sudo dnf config-manager --add-repo https://download.docker.com/linux/fedora/docker-ce.repo
sudo dnf install docker-ce docker-ce-cli containerd.io docker-compose-plugin
# Habilitar e iniciar el demonio systemd de Docker
sudo systemctl enable --now docker
# Adicionar el usuario actual al grupo docker
sudo usermod -aG docker $USER
\end{lstlisting}

\subsubsection{Debian}
Para distribuciones basadas en Debian (incluyendo Ubuntu y derivadas), el procedimiento estándar requiere la configuración del repositorio oficial mediante APT:
\begin{lstlisting}[language=bash]
# Instalacion de dependencias de red y criptografia
sudo apt-get update
sudo apt-get install ca-certificates curl
sudo install -m 0755 -d /etc/apt/keyrings
sudo curl -fsSL https://download.docker.com/linux/debian/gpg -o /etc/apt/keyrings/docker.asc
sudo chmod a+r /etc/apt/keyrings/docker.asc

# Configuracion del repositorio en las fuentes de APT
echo \
  "deb [arch=$(dpkg --print-architecture) signed-by=/etc/apt/keyrings/docker.asc] https://download.docker.com/linux/debian \
  $(. /etc/os-release && echo "$VERSION_CODENAME") stable" | \
  sudo tee /etc/apt/sources.list.d/docker.list > /dev/null

# Instalacion de Docker Engine, CLI y Compose Plugin
sudo apt-get update
sudo apt-get install docker-ce docker-ce-cli containerd.io docker-compose-plugin
# Adicionar el usuario actual al grupo docker
sudo usermod -aG docker $USER
\end{lstlisting}

\subsection{Entornos Windows}

En sistemas operativos Windows, el despliegue del motor de Docker puede implementarse mediante dos arquitecturas distintas:

\subsubsection{Arquitectura A: Docker Desktop (Implementación Híbrida/Gráfica)}
\begin{enumerate}
    \item Descargue el instalador oficial de \href{https://www.docker.com/products/docker-desktop}{Docker Desktop}.
    \item Durante el proceso de instalación, es crítico verificar que la integración con \textbf{WSL 2 (Windows Subsystem for Linux)} esté seleccionada, descartando el uso del hypervisor Hyper-V heredado.
    \item Inicialice la aplicación Docker Desktop. Este proceso arranca la máquina virtual ligera subyacente e inyecta los binarios de Docker en el PATH del sistema (PowerShell/CMD).
\end{enumerate}

\subsubsection{Arquitectura B: Docker sobre WSL 2 Nativo (Implementación CLI-only)}
Para minimizar el overhead del entorno gráfico y aislar el desarrollo a nivel de kernel de Linux, se recomienda instalar el motor de Docker directamente en una instancia de WSL 2:
\begin{enumerate}
    \item Despliegue el subsistema WSL junto a una distribución compatible (ej. Debian) ejecutando la siguiente instrucción en una instancia de PowerShell con privilegios de administrador:
\begin{lstlisting}[language=bash]
wsl --install -d Debian
\end{lstlisting}
    \item Acceda a la shell interactiva de la distribución desplegada.
    \item Ejecute las instrucciones de instalación detalladas en la sección correspondiente a \textbf{Debian} (ver arriba).
    \item \textit{Nota de compatibilidad WSL:} Las iteraciones recientes de WSL 2 integran soporte para \texttt{systemd}, permitiendo la gestión estándar de servicios. En arquitecturas no compatibles, la inicialización del demonio dockerd debe realizarse de forma explícita o mediante scripts de inicio (ej. \texttt{sudo service docker start}).
\end{enumerate}

\subsection{Validación del Entorno}
Independientemente de la arquitectura y sistema operativo anfitrión, ejecute las siguientes instrucciones en la terminal de preferencia para confirmar la accesibilidad y correcta inicialización de los binarios:
\begin{lstlisting}[language=bash]
docker --version
docker compose version
\end{lstlisting}
Una salida estándar informando las versiones de compilación correspondientes, sin retornos de error asociados a privilegios o descriptores de sockets inactivos, indicará que el entorno está operativo para el desarrollo.

\vspace{0.5cm}
\noindent\textbf{\textwarning\ Atención (Error: \texttt{failed to connect to the docker API /var/run/docker.sock}):} \\
Si al ejecutar comandos de Docker recibe un error indicando que no se puede conectar al socket, significa que el servicio/demonio de Docker no está en ejecución. 
\begin{itemize}
    \item En \textbf{Linux (Arch/Fedora/Debian):} Asegúrese de haber iniciado el servicio con \texttt{sudo systemctl start docker} (y habilitarlo con \texttt{enable --now}).
    \item En \textbf{Windows:} Asegúrese de tener abierta la aplicación de \textit{Docker Desktop}.
\end{itemize}

\vspace{0.5cm}
\noindent\textbf{\textwarning\ Atención (Error: \texttt{permission denied while trying to connect to the docker API}):} \\
Si recibe un error de permisos en Linux, significa que el demonio está corriendo pero su usuario no tiene privilegios para interactuar con él.
\begin{itemize}
    \item \textbf{Solución permanente (Recomendada):} Agregue su usuario al grupo docker y recargue los grupos:
\begin{lstlisting}[language=bash]
sudo usermod -aG docker $USER
newgrp docker
\end{lstlisting}
    \item \textbf{Solución rápida temporal:} Ejecute Docker con privilegios de administrador usando \texttt{sudo} (ej. \texttt{sudo docker compose up --build}).
\end{itemize}

\section{Levantar el Proyecto (Modo Desarrollo)}

\begin{enumerate}
    \item \textbf{Configurar las variables de entorno:} Copie el archivo de plantilla \texttt{.env.example} y renómbrelo a \texttt{.env}. Esto evitará advertencias de variables vacías.
\begin{lstlisting}[language=bash]
cp .env.example .env
\end{lstlisting}
    \item Abra su terminal y navegue hasta la carpeta raíz del proyecto (donde se encuentra el archivo \texttt{docker-compose.yml}).
    \item Ejecute el siguiente comando para construir las imágenes y levantar todo el ecosistema en segundo plano:
\begin{lstlisting}[language=bash]
docker compose up -d --build
\end{lstlisting}
    \textit{Nota: La primera vez tomará varios minutos mientras descarga las imágenes base (Alpine/Slim) y compila las dependencias de los microservicios.}

    \item \textbf{Verificar que todo esté corriendo:}
    Ejecute:
\begin{lstlisting}[language=bash]
docker compose ps
\end{lstlisting}
    Debería ver los contenedores (\texttt{api\_gateway}, \texttt{frontend}, \texttt{scraper}, etc.) con estado \texttt{Up}.
\end{enumerate}

\section{Comandos Útiles del Día a Día}

\begin{itemize}
    \item \textbf{Ver los logs (errores y prints):} \\
    Para ver qué está pasando en todos los servicios (útil para debugear):
\begin{lstlisting}[language=bash]
docker compose logs -f
\end{lstlisting}
    Si solo desea ver los errores del scraper:
\begin{lstlisting}[language=bash]
docker compose logs -f scraper_supermercados
\end{lstlisting}

    \item \textbf{Apagar el proyecto:} \\
    Cuando deje de trabajar, apague todo para liberar RAM en su PC:
\begin{lstlisting}[language=bash]
docker compose down
\end{lstlisting}

    \item \textbf{Reconstruir tras instalar una nueva librería:} \\
    Si añadió un nuevo paquete al \texttt{package.json} o \texttt{requirements.txt}, debe reconstruir la imagen:
\begin{lstlisting}[language=bash]
docker compose up -d --build
\end{lstlisting}
\end{itemize}

\section{Buenas Prácticas del Equipo}

\begin{enumerate}
    \item \textbf{Nunca subas binarios pesados al repositorio:} Para eso está el \texttt{.dockerignore}.
    \item \textbf{Usa variables de entorno:} No quemes contraseñas en el código, usa el archivo \texttt{.env} que se pasará automáticamente a los contenedores.
    \item \textbf{Pide ayuda si hay error de memoria:} Si Docker se "congela", avisa. Hemos limitado la RAM con \texttt{cgroups} en el \texttt{docker-compose.yml} para simular las restricciones del servidor Lenovo de producción.
\end{enumerate}

\end{document}
